\documentclass{article}
\usepackage{graphicx}
\usepackage{algorithm}
\usepackage{algpseudocode}
\usepackage{hyperref}
\usepackage{listings}

\title{Analysis of Adaptive Radix Tree (ART) Design and Range Query Implementation}
\author{Muyan Zhou}
\date{\today}

\begin{document}

\maketitle

\section{Design Logic of Adaptive Radix Tree (ART)}

\subsection{Core Design Principles and Motivation}
The Adaptive Radix Tree (ART) is a high-performance in-memory indexing structure that addresses the limitations of traditional radix trees and Tries. Its core design principles include:

\begin{itemize}
\item \textbf{Adaptive Node Sizes}: ART uses four different node types (\texttt{Node4}, \texttt{Node16}, \texttt{Node48}, \texttt{Node256}) that automatically adapt to the number of children, as shown in the \texttt{ArtNodeType} enum in \texttt{art.cpp}. By using different node types, ART minimizes memory overhead while maintaining fast access times.

\item \textbf{Path Compression}: ART employs prefix compression to reduce tree height, implemented via the \texttt{prefix} and \texttt{prefixLength} fields in the \texttt{ArtNode} structure. Also, nodes with only one child are compressed, reducing tree height and improving lookup times.

\end{itemize}

\subsection{Improvements Over Traditional Structures}
ART improves upon traditional radix trees and Tries in several ways:

\begin{itemize}
\item \textbf{Memory Efficiency}: Unlike fixed-size nodes in conventional radix trees, ART's adaptive nodes reduce memory consumption for sparse nodes. ART automatically upgrades/downgrades node types during insertion/deletion (implemented in \texttt{insertNode*} and \texttt{eraseNode*} functions).

\item \textbf{Time Efficiency}: ART's adaptive node types and path compression lead to faster lookups and insertions. For example, the \texttt{lookup} function uses SIMD operations for \texttt{Node16} and direct indexing for \texttt{Node256}, significantly speeding up key comparisons.

\end{itemize}

\subsection{Efficient Operations}
ART supports efficient operations through:

\begin{itemize}
\item \textbf{Insertion}: The \texttt{insert} function handles prefix splitting and node type conversion. It uses path compression and maintains sorted keys for efficient child lookup.

\item \textbf{Lookup}: The \texttt{lookup} function leverages prefix compression and optimized child search methods (e.g. SIMD in Node16, direct indexing in Node256).

\item \textbf{Deletion}: The \texttt{erase} function handles node consolidation and type downgrading when child counts fall below thresholds.
\end{itemize}

\section{Implementation of Range Query}
At the beginning, the range is made inclusive to simplify the implementation. 
The range query implementation in \texttt{rangeScan} employs an optimized depth-first traversal with the following key aspects:

\begin{enumerate}
\item \textbf{Prefix Accumulation}: Maintains a running prefix during traversal that represents the path from the root to the current node.

\item \textbf{Three-State Range Checking}:
\begin{itemize}
\item \texttt{equalLow}: True if current path exactly matches lower bound prefix
\item \texttt{equalHigh}: True if current path exactly matches upper bound prefix
\item \texttt{in\_range}: Integer flag indicating position relative to range (1=inside, 0=border, -1=outside)
\end{itemize}

\item \textbf{Early Termination}:
\begin{itemize}
\item Skips subtrees when \texttt{in\_range} is -1
\item Fast-tracks collection when \texttt{in\_range} is 1
\end{itemize}

\item \textbf{Leaf Validation}:
\begin{itemize}
\item For leaves, performs exact key comparison against bounds
\item Only adds to results if key $\in$ [lowerKey, upperKey]
\end{itemize}
\end{enumerate}

\end{document}